\documentclass[11pt,a4paper]{ctexart}

\usepackage[a4paper,left=22mm,right=22mm,top=20mm,bottom=22mm]{geometry}
\usepackage{amsmath,amssymb}
\usepackage{booktabs,longtable,array}
\usepackage{enumitem}
\usepackage{xcolor}
\usepackage{hyperref}
\usepackage{listings}
\usepackage{fancyhdr}

\hypersetup{
  colorlinks=true,
  linkcolor=blue!55!black,
  urlcolor=blue!55!black,
  citecolor=blue!55!black,
  pdftitle={CCUSDT 数据结构传递小册子 v0.1},
  pdfauthor={Jiang / Codex}
}

\setlength{\parindent}{2em}
\setlength{\parskip}{0.35em}
\setlength{\headheight}{14pt}
\linespread{1.12}
\pagestyle{fancy}
\fancyhf{}
\fancyhead[L]{CCUSDT 数据结构传递小册子}
\fancyhead[R]{v0.1}
\fancyfoot[C]{\thepage}

\definecolor{codebg}{RGB}{247,248,247}
\definecolor{rulegray}{RGB}{210,215,210}
\definecolor{deepgreen}{RGB}{18,88,62}
\definecolor{softblue}{RGB}{36,84,128}

\lstdefinestyle{code}{
  basicstyle=\ttfamily\small,
  backgroundcolor=\color{codebg},
  frame=single,
  rulecolor=\color{rulegray},
  breaklines=true,
  columns=fullflexible,
  keepspaces=true,
  showstringspaces=false,
  tabsize=2
}
\lstset{style=code}

\newcommand{\code}[1]{\texttt{\detokenize{#1}}}
\newcommand{\term}[1]{\textbf{#1}}
\newcommand{\file}[1]{\texttt{\detokenize{#1}}}
\newcommand{\ds}{\displaystyle}

\title{\Huge CCUSDT 数据结构传递小册子\\[6pt]
\large 从市场事件到 Monitor 因果链的手动推导}
\author{面向未来自己与新 Codex}
\date{v0.1 / 2026-06-01}

\begin{document}
\maketitle

\begin{center}
\begin{minipage}{0.82\textwidth}
本册不是系统白皮书，也不是 API 手册。它只解决一个具体问题：
\begin{center}
\term{一条市场事实，为什么会在 Runner、Bot、事件日志和 Monitor 页面中变成不同的数据结构？}
\end{center}
读完以后，读者应该能拿起一条 \code{events.ndjson} 记录，手动执行
\code{pickIntentId}、\code{ensureChain}、\code{mergeSignal}、\code{mergeCapacity}
和 \code{updateChainFromEvent}，最后填出 Monitor 表格中的一行。
\end{minipage}
\end{center}

\newpage
\tableofcontents
\newpage

\section{第 0 章：为什么数据结构会不断变形}

\subsection{本章要解决的困惑}

初看 \file{systems/ccusdt_replay_exchange} 时，最容易困惑的地方不是代码语言，而是同一件事在不同位置叫法不同。

例如，策略想做空 CCUSDT。人在口头上会说：

\begin{quote}
这里有一个 short entry。
\end{quote}

但系统不会把这句话原样传到底。它会变成一串不同对象：

\begin{lstlisting}
market event
-> online feature state
-> decision frame
-> shadow signal
-> capacity decision
-> order intent
-> scheduled order
-> arrived order
-> ack / reject
-> fill
-> portfolio state
-> compact event log
-> monitor chain
\end{lstlisting}

这不是工程上的啰嗦，而是信息边界。每一层都只回答它应该回答的问题。

\subsection{第一性原理：结构变形来自所有权变化}

把系统拆成三类所有权：

\begin{longtable}{p{0.18\textwidth}p{0.35\textwidth}p{0.37\textwidth}}
\toprule
所有者 & 它拥有的事实 & 它不能越界做的事 \\
\midrule
Runner & 市场真相、虚拟时钟、订单到达、成交、账户 & 不能读取研究 label，不能决定策略因子 \\
Bot & 在线特征、策略状态、下单意图 & 不能控制 replay clock，不能读取 \code{date/} 或未来 PnL/MFE/MAE \\
Monitor & 事件日志和只读状态的可视化解释 & 不能创造事实，不能下单或撤单 \\
\bottomrule
\end{longtable}

用一句更硬的工程边界说：Runner 不读 label；Bot 不读 \code{date/}；Monitor 不创造事实。任何把这三条边界打穿的实现，都会让数据结构看起来更方便，但会让回测失去可审计性。

因此，数据结构的变形可以写成：

\[
\text{Fact}
\xrightarrow{\text{state update}}
\text{State}
\xrightarrow{\text{policy}}
\text{Intent}
\xrightarrow{\text{exchange simulation}}
\text{Execution}
\xrightarrow{\text{audit}}
\text{Chain}.
\]

这里的每一个箭头都是一个边界。边界越清楚，回测越不容易偷看未来。

\subsection{贯穿全书的真实样例}

本册采用已有 run：

\begin{lstlisting}
systems/ccusdt_replay_exchange/runs/det_barrier_full_span60s_20260521
\end{lstlisting}

这个 run 的关键摘要：

\begin{lstlisting}
run_id: det_barrier_full_span60s_20260521
source_label: canonical_sparse_exchange_stream_v1:CCUSDT:2026-05-18
determinism_mode: deterministic_replay_v1
public_stream_mode: batched_public_barrier_v1
fill_model: top_of_book_taker_ioc_v1
orders_submitted: 606
fills_created: 606
arrival_quote_lag_count: 0
\end{lstlisting}

后面会取最早的一组 intent：

\begin{lstlisting}
intent_id = 1  entry sell
intent_id = 2  exit buy
\end{lstlisting}

它们构成一笔 short round trip。我们会手动演示日志如何变成 Monitor 表格行。

\section{第 1 章：全链路数据结构地图}

\subsection{一张总图}

用文本图表示全链路：

\begin{lstlisting}
canonical market truth
  market_quote / market_trade / market_l2_update
        |
        v
Bot online state
  feature_snapshot
  decision_frame
        |
        v
strategy decision
  shadow_signal
  capacity_decision
        |
        v
order ingress
  submit_orders / order_intent
        |
        v
Runner execution
  order_scheduled
  order_arrived
  order_ack / order_reject
  fill
  portfolio_state_on_change
        |
        v
compact event log
  events.ndjson
  summary.json
        |
        v
Monitor read model
  Chain[intent_id]
  eventTypeCounts
  warnings
\end{lstlisting}

\subsection{每个结构回答的问题}

\begin{longtable}{p{0.22\textwidth}p{0.36\textwidth}p{0.32\textwidth}}
\toprule
数据结构 & 回答什么问题 & 不应该包含什么 \\
\midrule
\code{market_quote} & 交易所当前最优买卖价是什么 & 策略判断、未来收益 \\
\code{feature_snapshot} & Bot 截止当前能安全知道什么 & future label、\code{date/} 结果 \\
\code{shadow_signal} & 策略为什么想 entry 或 exit & 真实成交价 \\
\code{capacity_decision} & 杠杆约束下实际能给多少仓位 & 未来是否盈利 \\
\code{order_intent} & Bot 想向交易所提交什么订单 & 成交结果 \\
\code{order_arrived} & 订单到达时盘口是什么 & 策略重新解释 \\
\code{fill} & 交易所模拟成交了什么 & 未发生的订单 \\
\code{portfolio_state} & 成交后账户状态如何 & 研究标签 \\
\code{monitor chain} & 这笔订单的因果链是什么 & 新事实、新成交 \\
\bottomrule
\end{longtable}

\subsection{变换公式}

事件日志进入 Monitor 后，核心不是“读 JSON”，而是一个聚合变换：

\[
E_i \longmapsto k_i=\operatorname{pickIntentId}(E_i)
\]

\[
C_{k_i}^{(i+1)}
=\operatorname{updateChainFromEvent}\bigl(C_{k_i}^{(i)}, E_i\bigr).
\]

其中 \(E_i\) 是第 \(i\) 条 \code{events.ndjson} 事件，\(k_i\) 是 \code{intent_id}，\(C_k\) 是 Monitor 中按 \code{intent_id} 聚合出的订单链。

如果一条事件没有 \code{intent_id}，它仍然可能影响总统计，例如 \code{eventTypeCounts}，但不会进入某条订单链。

\section{第 2 章：市场事件到 Bot 在线状态}

\subsection{public stream 的最小含义}

Runner 对 Bot 推送的 public stream 只代表交易所公开可见信息。常见事件包括：

\begin{lstlisting}
market_quote
market_trade
market_l2_update
market_decision_frame
\end{lstlisting}

这里最重要的公共字段是：

\begin{longtable}{p{0.24\textwidth}p{0.64\textwidth}}
\toprule
字段 & 含义 \\
\midrule
\code{seq} & Runner 内部公共事件序号，常用于重新定位观察点 \\
\code{local_ts_us} & 本地接收或 canonical 记录的微秒时间戳 \\
\code{quote_seq} & quote frame 序号 \\
\code{payload.quote} & bid、ask、mid 等盘口一档字段 \\
\bottomrule
\end{longtable}

这些字段只说“市场发生了什么”。它们还不是策略。

\subsection{从事件到状态}

Bot 不应该对每条事件都直接下判断。它先维护在线状态：

\[
S_t = U(S_{t-1}, E_t),
\]

其中 \(E_t\) 是当前市场事件，\(S_t\) 是 Bot 截止当前时刻的状态。

在当前策略里，\(S_t\) 至少包含：

\begin{itemize}[leftmargin=2em]
\item quote 状态：bid、ask、mid、spread、\code{frames_since_mid_change}；
\item trade flow 状态：rolling TFI、成交数量、买卖方向；
\item L2 摘要状态：bid/ask levels、best bid/ask；
\item private state：当前仓位、pending order、订单数量；
\item closed outcome state：只由已经关闭的历史 entry 生成 R5。
\end{itemize}

\subsection{feature snapshot 回答的问题}

\code{feature_snapshot} 是 Bot 把在线状态打包成一个可审计对象。它回答：

\begin{quote}
在 \code{observed_seq} 这个观察点，策略能安全知道什么？
\end{quote}

一个缩略例子：

\begin{lstlisting}
{
  "schema_id": "ccusdt_online_feature_snapshot_v1",
  "runtime_safe": true,
  "observed_seq": 1144,
  "local_ts_us": 1779062636310942,
  "quote": {
    "bid": 0.15406,
    "ask": 0.15408,
    "mid": 0.15407,
    "spread_bps": 1.2981
  },
  "trade_flow": {
    "tfi": -1.0,
    "rolling_trades": 1
  },
  "closed_outcomes": {
    "r5_ready": true,
    "r5": 0.0
  }
}
\end{lstlisting}

\code{runtime_safe=true} 是本章的核心。它不是装饰字段，而是边界声明：这个 snapshot 不能使用未来收益、未来路径、\code{date/} 研究标签、scored entries、MFE、MAE 或 PnL label。

\section{第 3 章：Bot 状态到订单意图}

\subsection{从状态到策略信号}

策略不是直接输出订单。它先输出 \code{shadow_signal}。这个结构回答：

\begin{quote}
为什么此刻应该 entry 或 exit？
\end{quote}

对四象限策略来说，核心字段包括：

\begin{longtable}{p{0.28\textwidth}p{0.58\textwidth}}
\toprule
字段 & 含义 \\
\midrule
\code{event_kind} & \code{shadow_entry} 或 \code{shadow_lifecycle_close} \\
\code{cell} & 四象限单元，例如 \code{00_none}、\code{10_r5_only} \\
\code{side} & Bot 想提交的方向，\code{buy} 或 \code{sell} \\
\code{observed_seq} & 这个判断绑定的市场观察点 \\
\code{entry_ts_us} & entry 触发时间 \\
\code{requested_exposure} & 策略想要的敞口 \\
\code{actual_exposure} & 容量约束后实际敞口 \\
\bottomrule
\end{longtable}

\subsection{容量约束为什么是单独结构}

策略想要下单，不代表账户能承载这笔风险。于是有 \code{capacity_decision}：

\begin{lstlisting}
{
  "schema_id": "ccusdt_capacity_decision_v1",
  "capacity_profile": "core_idle01",
  "capacity_source": "core",
  "leverage_cap": 3.0,
  "requested_exposure": 0.625,
  "actual_exposure": 0.625,
  "clipped_exposure": 0.0,
  "open_exposure_before": 0.0,
  "open_exposure_after": 0.625,
  "shadow_position_id": 3372,
  "skipped": false
}
\end{lstlisting}

这个结构回答：

\begin{quote}
在当前打开仓位和 3x cap 下，这个信号实际能给多少仓位？
\end{quote}

如果没有这层结构，策略结果就会把 alpha、杠杆、并发、clipping 混在一起。

\subsection{从信号到订单意图}

Bot 最后输出的是 \code{submit_orders}。它不能决定成交价，只能告诉 Runner：

\begin{quote}
我在 \code{observed_seq} 看到信号，现在想提交一个订单。
\end{quote}

缩略结构：

\begin{lstlisting}
{
  "type": "submit_orders",
  "observed_seq": 1144,
  "observed_local_ts_us": 1779062636310942,
  "orders": [
    {
      "client_order_id": "shadow-entry-3372",
      "side": "sell",
      "kind": "market",
      "qty": 4.056729302567098,
      "tif": "ioc"
    }
  ]
}
\end{lstlisting}

注意这里没有 \code{fill_price}。成交价格属于 Runner。

\section{第 4 章：Runner 如何把意图变成成交}

\subsection{Runner 的职责}

Runner 是本地交易所。它拥有：

\begin{itemize}[leftmargin=2em]
\item virtual clock；
\item public/private stream；
\item order ingress；
\item latency；
\item fill model；
\item portfolio；
\item compact event log。
\end{itemize}

它不拥有策略 alpha，也不读取 label。

\subsection{order scheduled}

当 Runner 收到 intent 后，第一步是安排订单到达：

\begin{lstlisting}
{
  "event_type": "order_scheduled",
  "payload": {
    "intent_id": 1,
    "observed_seq": 1144,
    "observed_local_ts_us": 1779062636310942,
    "arrival_local_ts_us": 1779062636310942,
    "latency_us": 0,
    "arrival_mode": "timer"
  }
}
\end{lstlisting}

数学上：

\[
t_{\mathrm{arrival}}
= t_{\mathrm{observed}} + \ell,
\]

其中 \(\ell\) 是 latency profile 给出的延迟。样例里 \(\ell=0\)，所以 observed 和 arrival 的目标时间相同。

\subsection{order arrived}

\code{order_arrived} 是最关键的执行结构。它回答：

\begin{quote}
订单真正到达交易所时，盘口是什么？
\end{quote}

样例 intent 1：

\begin{lstlisting}
{
  "event_type": "order_arrived",
  "payload": {
    "intent_id": 1,
    "observed_seq": 1144,
    "arrival_seq": 1144,
    "arrival_quote_lag": 0,
    "observed_quote": {
      "bid": 0.15406,
      "ask": 0.15408,
      "mid": 0.15407,
      "seq": 535
    },
    "arrival_quote": {
      "bid": 0.15406,
      "ask": 0.15408,
      "mid": 0.15407,
      "seq": 535
    },
    "spread_bps_at_arrival": 1.2981,
    "latency_slippage_bps": 0.0,
    "order": {
      "client_order_id": "shadow-entry-3372",
      "side": "sell",
      "kind": "market",
      "qty": 4.056729302567098
    }
  }
}
\end{lstlisting}

这里要区分三件事：

\begin{longtable}{p{0.26\textwidth}p{0.58\textwidth}}
\toprule
对象 & 含义 \\
\midrule
\code{observed_quote} & Bot 生成 intent 时绑定的盘口 \\
\code{arrival_quote} & Runner 判定订单到达时使用的盘口 \\
\code{fill_price} & fill model 根据 arrival quote 算出的成交价 \\
\bottomrule
\end{longtable}

\subsection{top-of-book taker IOC 成交公式}

当前默认模型是 \code{top_of_book_taker_ioc_v1}：

\[
\operatorname{fill\_price} =
\begin{cases}
\operatorname{arrival.ask}, & \text{market buy},\\
\operatorname{arrival.bid}, & \text{market sell}.
\end{cases}
\]

因此 intent 1 是 sell，成交价是 arrival bid：

\[
\operatorname{fill\_price}_1 = 0.15406.
\]

intent 2 是 buy，成交价是 arrival ask：

\[
\operatorname{fill\_price}_2 = 0.15397.
\]

这也解释了一个常见问题：\code{fee_bps=0} 不等于零执行成本。主动买卖会 crossing spread；手续费和 spread cost 是两类成本。

\subsection{fill 与 portfolio state on change}

\code{fill} 记录成交事实：

\begin{lstlisting}
{
  "event_type": "fill",
  "payload": {
    "payload": {
      "intent_id": 1,
      "execution_model": "top_of_book_taker_ioc_v1",
      "fill_price": 0.15406,
      "fee_bps": 0.0,
      "fill": {
        "side": "sell",
        "price": 0.15406,
        "qty": 4.056729302567098,
        "liquidity": "taker"
      }
    }
  }
}
\end{lstlisting}

portfolio 事件则回答：

\begin{quote}
这次成交如何改变账户和 position lot？
\end{quote}

Monitor 可以展示 portfolio，但不能重新计算真实账户状态。账户状态仍以 Runner 事件为准。

\section{第 5 章：事件日志为什么是 NDJSON}

\subsection{一行一个事件}

\code{events.ndjson} 的含义是 newline-delimited JSON：每一行是一条完整 JSON 事件。

它适合 replay exchange 的原因：

\begin{itemize}[leftmargin=2em]
\item 可以边跑边写，不需要等整个 run 结束；
\item 可以 streaming parse，不必一次性读入几十 MB 或几百 MB；
\item 每一行都是独立审计单位；
\item compact log 可以只保存因果关键事件。
\end{itemize}

\subsection{compact log 不记录每帧 hold}

Bot 在大多数时刻不下单。如果每帧都写一条 hold，日志会膨胀成市场事件数量级：

\[
O(\text{market events}).
\]

compact log 只保留：

\[
O(\text{orders} + \text{fills} + \text{state changes} + \text{summary}).
\]

这不是丢信息，而是把“无动作状态机内部事实”和“外部可审计因果事件”分开。Monitor 第一版也沿用这个原则：它不展示每个 hold，只展示能解释订单的因果链。

\subsection{为什么按 intent id 审计}

一个 order lifecycle 会跨多条事件：

\begin{lstlisting}
capacity_decision
order_intent
order_scheduled
order_arrived
order_ack
fill
position_opened / position_closed
portfolio_state_on_change
\end{lstlisting}

这些事件可能在日志中相邻，也可能夹杂着 public batch barrier、feature checkpoint、progress 等事件。Monitor 不能靠行号理解订单，只能靠 \code{intent_id} 聚合。

聚合公式是：

\[
\mathcal{C}_k = \{E_i: \operatorname{pickIntentId}(E_i)=k\}.
\]

\section{第 6 章：手动执行 monitor parser}

本章对应当前文件：

\begin{lstlisting}
systems/ccusdt_replay_exchange/monitor/src/monitor_server.mjs
\end{lstlisting}

目标是从代码理解 Monitor 的读模型，而不是泛泛解释 JavaScript。

\subsection{pickIntentId：先找到订单链的 key}

函数逻辑可写成：

\begin{lstlisting}
function pickIntentId(event) {
  const p = event.payload ?? {};
  const inner = p.payload ?? {};
  const intent = p.intent ?? inner.intent ?? {};
  return p.intent_id ?? inner.intent_id ?? intent.intent_id
      ?? p.order?.intent_id ?? inner.order?.intent_id ?? null;
}
\end{lstlisting}

为什么要找这么多位置？因为事件结构有嵌套：

\begin{itemize}[leftmargin=2em]
\item \code{order_scheduled} 的 \code{intent_id} 在 \code{event.payload.intent_id}；
\item \code{fill} 的核心字段在 \code{event.payload.payload.intent_id}；
\item 某些策略响应可能把 intent 放在 \code{payload.intent} 里。
\end{itemize}

手算检查点：

\begin{lstlisting}
E = order_arrived event for intent 1
p = E.payload
p.intent_id = 1
pickIntentId(E) = 1
\end{lstlisting}

\subsection{ensureChain：没有链就创建空链}

Monitor 的 read model 是一个 map：

\[
\operatorname{chains}: \operatorname{intent\_id} \mapsto \operatorname{Chain}.
\]

如果第一次看到 \code{intent_id=1}，就创建：

\begin{lstlisting}
{
  intent_id: 1,
  event_ids: [],
  events: [],
  status: "pending",
  side: null,
  cell: null,
  capacity_source: null,
  requested_exposure: null,
  actual_exposure: null,
  observed_seq: null,
  arrival_seq: null,
  observed_quote: null,
  arrival_quote: null,
  fill_price: null,
  raw: {}
}
\end{lstlisting}

这个空链就是 Monitor 表格行的初始状态。

\subsection{mergeSignal：把策略解释合进链}

\code{shadow_signal} 不是成交事实，它是策略解释。Monitor 读取它来填：

\begin{lstlisting}
side
cell
shadow_position_id
observed_seq
requested_exposure
actual_exposure
capacity_source
\end{lstlisting}

对应规则：

\[
C.\operatorname{side} \leftarrow
C.\operatorname{side}\ \text{if exists, otherwise}\ signal.\operatorname{side}.
\]

代码里用的是 JavaScript 的空值合并思想：

\begin{lstlisting}
chain.side = chain.side ?? signal.side ?? signal.direction_label ?? null;
\end{lstlisting}

含义是：已有值优先，不用后来的事件随意覆盖前面已经确定的事实。

\subsection{mergeCapacity：把容量决策合进链}

\code{capacity_decision} 填的是容量字段：

\begin{lstlisting}
capacity_source
requested_exposure
actual_exposure
clipped_exposure
shadow_position_id
skipped
\end{lstlisting}

手算 intent 1：

\begin{lstlisting}
capacity_source = "core"
requested_exposure = 0.625
actual_exposure = 0.625
clipped_exposure = 0.0
skipped = false
\end{lstlisting}

所以 Monitor 表格里 intent 1 的 capacity 列是 \code{core}，requested 和 actual 都是 \(0.625\)。

\subsection{updateChainFromEvent：按事件类型更新链}

主函数的结构是：

\begin{lstlisting}
intentId = pickIntentId(event)
chain = ensureChain(chains, intentId)
mergeSignal(chain, signal)
mergeCapacity(chain, capacity)

if event_type == "order_scheduled":
    fill arrival_local_ts_us and bridge latency
elif event_type == "order_arrived":
    fill arrival_seq, observed_quote, arrival_quote, spread
elif event_type == "order_ack":
    fill status, order id, fill price if present
elif event_type == "fill":
    status = "filled"; fill fill_price and fill_qty
elif event_type == "portfolio_state_on_change":
    attach account
\end{lstlisting}

这不是创造新信息，只是把不同事件里的字段投影到同一个 \code{Chain} 对象上。

\subsection{parseEvents：streaming 读取大日志}

\code{events.ndjson} 可能很大。Monitor 不应该：

\begin{lstlisting}
readFileSync(events.ndjson)
JSON.parse(wholeFile)
\end{lstlisting}

而是逐行：

\begin{lstlisting}
for each line:
    event = JSON.parse(line)
    eventTypeCounts[event.event_type] += 1
    updateChainFromEvent(chains, event)
\end{lstlisting}

最终输出：

\begin{lstlisting}
{
  event_count,
  bad_json,
  chains: sorted Chain rows,
  eventTypeCounts,
  has_events
}
\end{lstlisting}

\subsection{normalizeSummary：把 run 摘要变成页面头部}

\code{summary.json} 和 \code{profile_manifest.json} 的字段不完全重复。Monitor 需要把它们合成页面头部的统一对象：

\begin{lstlisting}
transport_profile = summary.public_stream_mode
                 ?? profileManifest.transport_profile.name
fill_profile = summary.fill_model
            ?? profileManifest.fill_profile.name
orders_submitted = summary.orders_submitted
fills_created = summary.fills_created
arrival_quote_lag_count = summary.arrival_quote_lag_count
\end{lstlisting}

这一步回答的是页面问题：

\begin{quote}
这个 run 是什么 profile？订单和成交是否对齐？是否有 arrival quote lag？
\end{quote}

\section{第 7 章：一个真实订单链的手算案例}

\subsection{样例目标}

我们手动构造 intent 1 的 Monitor 行。真实 run 中 intent 1 是 entry sell：

\begin{lstlisting}
client_order_id = shadow-entry-3372
side = sell
cell = 00_none
capacity_source = core
requested_exposure = 0.625
actual_exposure = 0.625
observed_seq = 1144
arrival_seq = 1144
fill_price = 0.15406
spread_bps_at_arrival = 1.2981
latency_slippage_bps = 0.0
\end{lstlisting}

\subsection{第 1 步：读 capacity decision}

事件 4 的关键字段：

\begin{lstlisting}
event_type = "capacity_decision"
intent_id = 1
capacity_decision.capacity_source = "core"
capacity_decision.requested_exposure = 0.625
capacity_decision.actual_exposure = 0.625
capacity_decision.clipped_exposure = 0.0
shadow_signal.cell = "00_none"
shadow_signal.side = "sell"
shadow_signal.observed_seq = 1144
\end{lstlisting}

手动执行：

\begin{lstlisting}
k = pickIntentId(E4) = 1
C1 = ensureChain(chains, 1)
mergeSignal(C1, E4.payload.shadow_signal)
mergeCapacity(C1, E4.payload.capacity_decision)
\end{lstlisting}

此时：

\begin{lstlisting}
C1.status = "pending"
C1.side = "sell"
C1.cell = "00_none"
C1.capacity_source = "core"
C1.requested_exposure = 0.625
C1.actual_exposure = 0.625
C1.observed_seq = 1144
\end{lstlisting}

\subsection{第 2 步：读 order scheduled}

事件 6：

\begin{lstlisting}
event_type = "order_scheduled"
payload.intent_id = 1
payload.observed_seq = 1144
payload.observed_local_ts_us = 1779062636310942
payload.arrival_local_ts_us = 1779062636310942
payload.latency_us = 0
\end{lstlisting}

执行后：

\begin{lstlisting}
C1.events += ["order_scheduled"]
C1.arrival_local_ts_us = 1779062636310942
C1.status = "scheduled"
\end{lstlisting}

这里没有成交价，因为订单还没有到达交易所。

\subsection{第 3 步：读 order arrived}

事件 8：

\begin{lstlisting}
event_type = "order_arrived"
payload.intent_id = 1
payload.arrival_seq = 1144
payload.observed_quote.bid = 0.15406
payload.observed_quote.ask = 0.15408
payload.arrival_quote.bid = 0.15406
payload.arrival_quote.ask = 0.15408
payload.spread_bps_at_arrival = 1.2981
payload.latency_slippage_bps = 0.0
\end{lstlisting}

执行后：

\begin{lstlisting}
C1.arrival_seq = 1144
C1.observed_quote = {bid: 0.15406, ask: 0.15408, mid: 0.15407}
C1.arrival_quote = {bid: 0.15406, ask: 0.15408, mid: 0.15407}
C1.spread_bps_at_arrival = 1.2981
C1.latency_slippage_bps = 0.0
C1.status = "arrived"
\end{lstlisting}

因为 \code{arrival_quote_lag=0}，observed 和 arrival 使用同一盘口。真实有延迟时，这两个对象可能不同。

\subsection{第 4 步：读 fill}

事件 10：

\begin{lstlisting}
event_type = "fill"
payload.payload.intent_id = 1
payload.payload.fill_price = 0.15406
payload.payload.fill.side = "sell"
payload.payload.fill.qty = 4.056729302567098
payload.payload.fee_bps = 0.0
\end{lstlisting}

执行后：

\begin{lstlisting}
C1.status = "filled"
C1.fill_price = 0.15406
C1.fill_qty = 4.056729302567098
C1.side = "sell"
\end{lstlisting}

根据成交公式，sell market IOC 打 bid：

\[
\operatorname{fill\_price}_1
=\operatorname{arrival.bid}
=0.15406.
\]

\subsection{第 5 步：填 Monitor 表格}

最终 Monitor 行：

\begin{longtable}{p{0.22\textwidth}p{0.58\textwidth}}
\toprule
列 & 值 \\
\midrule
\code{intent} & 1 \\
\code{status} & \code{filled} \\
\code{side} & \code{sell} \\
\code{cell} & \code{00_none} \\
\code{capacity} & \code{core} \\
\code{requested} & 0.625 \\
\code{actual} & 0.625 \\
\code{obs seq} & 1144 \\
\code{arr seq} & 1144 \\
\code{fill} & 0.15406 \\
\code{spread} & 1.2981 bps \\
\code{lat slippage} & 0.0 bps \\
\bottomrule
\end{longtable}

这就是“手动执行函数”的意义：Monitor 页面不是独立计算系统，而是把日志里的分散事实按 \code{intent_id} 重新排成可读形式。

\section{第 8 章：常见混淆与调试路线}

\subsection{混淆一：nested payload}

很多事件长这样：

\begin{lstlisting}
event.payload.payload.intent_id
\end{lstlisting}

这不是错误，而是因为外层 payload 是 event log 包装，内层 payload 是 private stream 消息体。调试时不要只搜第一层 \code{payload.intent_id}。

\subsection{混淆二：observed quote 与 arrival quote}

\code{observed_quote} 是 Bot 决策绑定的盘口；\code{arrival_quote} 是 Runner 订单到达时用于成交的盘口。

如果 latency profile 是 deterministic 且 \(\ell=0\)，两者常常一样。如果引入真实延迟或 pressure mode，它们可能不同。

\[
\text{latency slippage}
\approx
\operatorname{value}(\text{arrival quote})
-
\operatorname{value}(\text{observed quote}).
\]

\subsection{混淆三：strategy intent 与 exchange order}

\code{order_intent} 是 Bot 的请求；\code{order_ack} 是交易所模拟器的回应；\code{fill} 是成交事实。

因此：

\begin{lstlisting}
strategy intent != exchange order != fill
\end{lstlisting}

如果一个 intent 被 risk reject，它不会自然生成 fill。如果一个 order partial fill，fill qty 也不应等同于 requested qty。

\subsection{混淆四：fee 与 spread}

当前 venue-fee baseline 可以是：

\[
\operatorname{fee\_bps}=0.
\]

但 taker crossing spread 仍然存在：

\[
\operatorname{long\ round\ trip}
=10^4\log\frac{b_{\mathrm{exit}}}{a_{\mathrm{entry}}}.
\]

这里 entry 用 ask，exit 用 bid。即使没有手续费，也可能因为 crossing spread 和路径衰减损失很多。

\subsection{混淆五：fast 与 strict}

fast backtest 回答：

\begin{quote}
策略、容量、尾部和参数大概如何？
\end{quote}

strict replay 回答：

\begin{quote}
在交易所式 public/private stream、latency、fill 和 event log 下，行为是否可信？
\end{quote}

Monitor 第一版主要服务 strict replay 的解释，因为它展示的是 Runner 事件日志。

\subsection{调试路线}

如果 Monitor 某一行看起来不对，按这个顺序查：

\begin{enumerate}[leftmargin=2em]
\item 在 \code{events.ndjson} 搜 \code{"intent_id":N}；
\item 确认是否有 \code{capacity_decision}；
\item 确认是否有 \code{order_scheduled}；
\item 确认是否有 \code{order_arrived}；
\item 确认是否有 \code{order_ack} 或 \code{order_reject}；
\item 确认是否有 \code{fill}；
\item 对照 \file{monitor_server.mjs} 的 \code{updateChainFromEvent} 分支。
\end{enumerate}

\section*{附录 A：字段词典}
\addcontentsline{toc}{section}{附录 A：字段词典}

\begin{longtable}{p{0.26\textwidth}p{0.58\textwidth}}
\toprule
字段 & 解释 \\
\midrule
\code{intent_id} & Bot 提交意图后，Runner/日志用于串联生命周期的 key \\
\code{observed_seq} & Bot 观察并触发 intent 的市场序号 \\
\code{arrival_seq} & Runner 判定订单到达时对应的市场/stream 序号 \\
\code{observed_local_ts_us} & 触发时刻的本地微秒时间 \\
\code{arrival_local_ts_us} & 订单目标到达时间 \\
\code{observed_quote} & 策略观察时的盘口 \\
\code{arrival_quote} & 成交计算使用的到达盘口 \\
\code{fill_price} & fill model 给出的成交价 \\
\code{spread_bps_at_arrival} & 到达盘口的 bid-ask spread，单位 bps \\
\code{latency_slippage_bps} & latency 导致的相对价格偏移 \\
\code{capacity_source} & core 或 idle sleeve 等容量来源 \\
\code{requested_exposure} & 策略请求敞口 \\
\code{actual_exposure} & 容量约束后的实际敞口 \\
\code{runtime_safe} & 字段只来自运行时可见信息，不包含未来 label \\
\bottomrule
\end{longtable}

\section*{附录 B：函数索引}
\addcontentsline{toc}{section}{附录 B：函数索引}

\begin{longtable}{p{0.30\textwidth}p{0.54\textwidth}}
\toprule
函数 & 作用 \\
\midrule
\code{pickIntentId} & 从不同嵌套位置取出 \code{intent_id} \\
\code{ensureChain} & 为某个 \code{intent_id} 创建或取回聚合链 \\
\code{mergeSignal} & 把 \code{shadow_signal} 的策略解释字段并入链 \\
\code{mergeCapacity} & 把 \code{capacity_decision} 的容量字段并入链 \\
\code{updateChainFromEvent} & 按事件类型更新链字段 \\
\code{parseEvents} & 逐行 streaming 读取 \code{events.ndjson} 并生成链表 \\
\code{normalizeSummary} & 合成页面头部 summary/profile 字段 \\
\bottomrule
\end{longtable}

\section*{附录 C：最小复现命令}
\addcontentsline{toc}{section}{附录 C：最小复现命令}

启动 Monitor 页面：

\begin{lstlisting}
npm --prefix systems/ccusdt_replay_exchange/monitor run dev -- --port 8810
\end{lstlisting}

检查离线 run API：

\begin{lstlisting}
Invoke-RestMethod `
  "http://127.0.0.1:8810/api/offline-run?runDir=det_barrier_full_span60s_20260521"
\end{lstlisting}

直接搜 intent：

\begin{lstlisting}
Select-String `
  -Path systems/ccusdt_replay_exchange/runs/det_barrier_full_span60s_20260521/events.ndjson `
  -Pattern '"intent_id":1'
\end{lstlisting}

编译本册：

\begin{lstlisting}
xelatex -interaction=nonstopmode `
  -output-directory docs/books/ccusdt-data-structure-booklet/out `
  docs/books/ccusdt-data-structure-booklet/book.tex
\end{lstlisting}

\end{document}
